\begin{textsample}{NEG Dim 8 – global\_north\_2024\_05 – Score -1.00 – global\_north\_2024\_05/108326385.txt}  \label{ex:f8_neg_381}
Israel Refutes South Africa ' s Accusations At UN World Court Law and Crime Prevention The UN International Court of Justice ( ICJ ), on Friday, heard the response from Israel on the case brought forward by South Africa requesting emergency provisional measures to immediately halt Israeli military operations under way in Rafah, in southern Gaza, where over one million Palestinians were sheltering after having been displaced from elsewhere in the enclave.
Linked to South Africa ' s ongoing case accusing Israel of violating its obligations under the Convention on the Prevention and Punishment of the Crime of Genocide ( the Genocide Convention ), the new request, filed on 10 May, asked the ICJ to order Israel to"immediately withdraw and cease its military operations in the Rafah governate".
Gilad Noam presents Israel ' s arguments at the ICJ.
An ' obscene exploitation ' Appearing before the Court, Gilad Noam, co-agent of Israel, refuted South Africa ' s claims, terming it an"obscene exploitation"of the"most sacred"Genocide Convention."South Africa presents @ @ @ @ @ @ @ @ @ @ scope of less than five months, with a picture that is completely divorced from the facts and circumstances."He stated that Israel is engaged in a"difficult and tragic"armed conflict, a fact essential to"comprehend the situation"but one that is ignored by South Africa."It makes a mockery of the heinous charge of genocide... facts matter and truth should matter.
Words must retain their meaning.
Calling something a genocide again and again does not make it genocide,"he added.
Israel did not start the war Advertisement - scroll to continue reading Are you getting our free newsletter?
Mr.
Noam further stated that it was not Israel that started the war, recalling the"horrific onslaught"on 7 October 2023 by Hamas and other Palestinian groups targeting Israeli civilians and communities, killing over 1,200 and taking 254 women, men and children hostage.
He added that Hamas and other terrorist groups in Gaza continue to attack Israel, displacing communities and @ @ @ @ @ @ @ @ @ @ use Palestinian civilians as human shields."Rafah, in particular, is a focal point for the ongoing terrorist activity,"he said, accusing Hamas of having"intricate underground tunnel infrastructure"with \textbf{command-and-control} rooms, military equipment, and potentially for smuggling Israeli hostages out of Gaza.
He also noted that even though the ICJ called for the immediate release of the hostages, they remain under captivity.
Not a large-scale operation"The reality is that any State put in Israel ' s difficult position would do the same.
The right of defence against the brutality of the Hamas terrorist organization can not be in doubt.
It is an inherent right afforded to Israel, as it is to any State,"Mr.
Noam said.
The Israeli representative stated his country ' s commitment to protecting itself,"in accordance with the law, which is why it has worked diligently to enable the protection of civilians, even as Hamas seeks deliberately to endanger them.""That is why @ @ @ @ @ @ @ @ @ @ but rather specific, limited and localized operations prefaced with evacuation efforts and support for humanitarian activities,"he added.
Fully and sincerely engaged Concluding his statement, Mr.
Noam cited the Court ' s rejection of South Africa ' s earlier requests for similar provisional measures, and added that it would be"wholly inappropriate"to grant a provisional measure under such terms."South Africa has not given sufficient reason why the Court should now deviate from or essentially duplicate its earlier decisions,"he said, noting that Israel is"engaged fully and sincerely"in the proceedings,"despite the outrageous and libelous claims levelled against it.""Israel has made clear time and again its unwavering commitment to its obligations under its international law.
It has done this while the fighting continues and its citizens are still under attack,"he said.
Did you know Scoop has an Ethical Paywall?
If you ' re using Scoop for work, your organisation needs to pay a small license fee @ @ @ @ @ @ @ @ @ @ because your organisation is benefiting from using our news resources.
In return, we ' ll also give your team access to pro news tools and keep Scoop free for personal use, because public access to news is important!

% matched lemmas: command-and-control
\end{textsample}
